\section{Discusión de decisiones de diseño}
\label{sec:decisiones}

Las decisiones de este proyecto están registradas una por una ---con su
contexto, las alternativas descartadas y la evidencia cuando la hubo--- en
una bitácora versionada junto al código~\cite{Bitacora2026}. Esta sección
discute tres de ellas, las que tenían una alternativa razonable que se
descartó con evidencia y no por preferencia.

\subsection{Primitivas: mutex y condiciones, no semáforos}

\textbf{Decisión.} Un mutex y cuatro variables de condición POSIX
configuradas para compartirse entre procesos. El enunciado permite las dos
vías y deja la elección al grupo, así que esta es la decisión de la que
cuelgan las dos siguientes.

\textbf{Alternativa descartada.} Semáforos POSIX compartidos entre procesos,
que es la construcción original de Dijkstra para este problema y necesita
menos código: dos contadores, un binario de exclusión y ningún ciclo de
predicado.

\textbf{Evidencia.} Tres razones, ninguna de preferencia.

\emph{Con semáforos hay un orden que interbloquea y hay que recordar no
escribirlo}, el de las dos \code{P} del consumidor
(sección~\ref{sec:lecturas}). Es justo lo que el enunciado prohíbe ---no
esperar mientras se conserva el bloqueo que necesita quien puede satisfacer
esa espera--- y con una variable de condición ese error no se puede
escribir.

\emph{El cierre no necesita permisos falsos.} Despertar a todos los
bloqueados es un \code{broadcast} sobre las cuatro condvars, y lo que cada
uno reevalúa al despertar es el estado real. Por la vía de los semáforos hay
que publicar permisos que no corresponden a ningún mensaje ni a ningún
espacio, y el propio enunciado pide entonces distinguirlos de los reales:
una obligación que esta decisión elimina en vez de cumplir.

\emph{La contabilidad queda en una sola variable.} Un \code{count} bajo un
mutex permite evaluar los dos predicados sobre el mismo estado, lo que con
dos semáforos no vale en todo instante (sección~\ref{sec:lecturas}).

\textbf{Lo que costó.} Las primitivas dejan de ser objetos del sistema y
pasan a vivir dentro del segmento, con \code{PTHREAD\_PROCESS\_SHARED} y un
reloj explícito en las condvars; alguien tiene que crearlas y alguien
destruirlas, que es de donde sale el protocolo de cierre con finalizador
único. Un semáforo con nombre lo habría resuelto el sistema operativo. Ese
es el precio que el proyecto paga por las tres razones de arriba.

\subsection{Tamaño del buffer: no hay un valor por omisión razonable}

\textbf{Decisión.} \code{--capacity} es obligatoria y no tiene valor por
omisión: \code{creator} se niega a crear la instancia si no se la dan.

\textbf{Alternativa descartada.} Darle un valor por omisión cómodo ---8,
digamos--- para poder invocar \code{creator} sin pensar en la capacidad, que
es lo que haría cualquiera que asuma que más espacio nunca estorba.

\textbf{Evidencia.} Ese valor cómodo no existe
(sección~\ref{sec:resultados}): subir la capacidad solo mejora el throughput
cuando ninguno de los dos lados es el cuello de botella (+18\,\% en 5/5~ms,
plano en 2/8~ms), mientras que la latencia sube siempre, hasta un factor de
66. En régimen saturado la latencia ni siquiera es una medición: es la
capacidad dividida por la tasa de drenaje. Cualquier valor por omisión
habría sido bueno para una de las dos métricas y malo para la otra, y el
programa no tiene forma de saber cuál le importa a quien lo corre.

\textbf{Consecuencia práctica.} Elegir la capacidad es elegir un punto en un
intercambio entre caudal y latencia, no una mejora gratuita, así que el
programa obliga a elegirlo de forma explícita. Una simulación de eventos
discretos del mismo escenario, sin mutex ni registro
(\code{scripts/modelo.py}), agrega que el beneficio satura en seco: de
capacidad~128 en adelante el throughput no gana nada, así que el 64 del
barrido ya estaba en el plateau y probar capacidades mayores no habría
aportado información.

\subsection{Granularidad del bloqueo: un solo mutex, y la evidencia de que alcanza}

\textbf{Decisión.} Un mutex por instancia, que protege el anillo, los
contadores, las tablas de participantes y las métricas. \code{buffer.h}
asume que ya está tomado.

\textbf{Alternativa descartada.} Bloqueo más fino ---un candado para el
anillo y otro para las métricas, o secciones críticas separadas por
campo--- con el argumento de reducir contención.

\textbf{Evidencia.} La contención no es el costo, así que el bloqueo fino
optimizaría lo que no duele. Adquirir el mutex cuesta menos de
1~\textmu s por operación en casi todas las celdas (0.24--0.63~\textmu s los
productores, 0.35--1.7~\textmu s los consumidores), entre tres y cuatro
órdenes de magnitud por debajo de las pausas y de las esperas por la
contraparte. Dividir el mutex habría multiplicado los estados posibles y los
órdenes de adquisición ---con el riesgo de interbloqueo que eso trae, el
mismo que Dijkstra señala para el orden de las operaciones--- a cambio de un
microsegundo.

Lo que sí es caro quedó identificado al analizar los resultados, y no es el
candado: despertar a un proceso dormido cuesta unos 60~\textmu s en este
entorno (\code{cond\_signal}, el despacho del planificador y la
readquisición del mutex), dos órdenes de magnitud más que tomar el mutex sin
contención. En el escenario 8/2~ms, donde el consumidor se bloquea en 197 de
cada 200 extracciones, eso es lo que domina la latencia media por mensaje.
La conclusión para el diseño es la inversa de la intuición: conviene
preocuparse por \emph{cuántas veces se duerme un proceso}, no por cuánto
dura la sección crítica.
